\documentclass[a4paper,11pt]{article}
\usepackage[T1]{fontenc}
\usepackage[utf8]{inputenc}
\usepackage{t1enc}
\usepackage[dvips]{color}
\usepackage{amsmath,amssymb,makeidx,stmaryrd}
\usepackage[normalem]{ulem}
\usepackage{fancybox}
\definecolor{gris}{gray}{0.9}
\usepackage{pst-plot,pst-tree,pstricks,pst-node,pst-text}
\usepackage{pst-eucl}
\usepackage{pstricks-add}
\usepackage[left=3.5cm, right=3.5cm, top=3cm, bottom=3cm]{geometry}
\newcommand{\vect}[1]{\overrightarrow{\,\mathstrut#1\,}}
\def\bas{\fcolorbox{gris}{gris}{\parbox{0.986\textwidth}{
	\color{red}\footnotesize 1993\hfill Page : \thepage/1}}}
\def\haut{\fcolorbox{gris}{gris}{\parbox{0.986\textwidth}{
	\color{red}
	\bf Concours Général de Mathématiques \hfill 1993
	}}}
\makeatletter
\def\ps@cg{%
\let\@oddhead\haut%
\let\@evenhead\haut%
\let\@oddfoot\bas
\let\@evenfoot\bas}
\makeatother
\newcommand{\R}{\mathbb{R}}
\newcommand{\N}{\mathbb{N}}
\newcommand{\D}{\mathbb{D}}
\newcommand{\Z}{\mathbb{Z}}
\newcommand{\Q}{\mathbb{Q}}
\newcommand{\C}{\mathbb{C}}
\pagestyle{cg}
\renewcommand{\theenumi}{\textbf{\arabic{enumi}}}
\renewcommand{\labelenumi}{\textbf{\theenumi.}}
\renewcommand{\theenumii}{\textbf{\alph{enumii}}}
\renewcommand{\labelenumii}{\textbf{\theenumii.}}
\parindent0pt
\everymath{\displaystyle}
\usepackage[french]{babel}

\DecimalMathComma
\usepackage[np]{numprint}
\begin{document}

\begin{center}\textbf{{\color{blue} Exercice I}}
\end{center}

On appelle \og boîte de poids \fg, tout ensemble de poids comprenant: 

\begin{center}
\begin{tabular}{l}
$x_1$ poids de $d_1$ grammes chacun\\
$x_2$ poids de $d_2$ grammes chacun\\
\dotfill\\
$x_k$ poids de $d_k$ grammes chacun\\
\end{tabular}
\end{center}

Les nombres $x_i$ et $d_i$ (pour $1 \leqslant i \leqslant k$) étant des entiers naturels non nuls, tels que :

\[1 \leqslant d_1 < d_2 < \dots  < d_k.\]

On pose $n=x_1d_1 + x_2d_2+ \ldots + x_kd_k$ et on dit que la masse totale de la boîte de poids est $n$ grammes.

Cette boîte de poids est dite \og parfaite\fg si elle permet d'obtenir, de manière unique, toute masse de $0, 1, \ldots , n$ grammes, c'est-à-dire si, pour tout entier $m$ de l'ensemble $\{0, 1, \ldots,~n\}$, il existe des entiers $y_1, y_2, \ldots , y_k$ uniques tels que pour tout indice $i$,\: $1 \leqslant i \leqslant k$, on ait :

$0 \leqslant y_i \leqslant x_i$ et $m = y_1d_1 + y_2d_2 + \ldots + y_kd_k$. 

\medskip

\begin{enumerate}
\item Déterminer toutes les boîtes de poids de masse totale $5$ grammes. Quelles sont celles qui sont parfaites ?
\item Montrer que pour une boîte de poids parfaite de masse totale $n$ grammes, on a, avec les notations du préambule:

\[\left(1 + x_1\right)\left(1 + x_2\right)\ldots \left(1 + x_k\right) = n + 1.\]

\item Inversement, on se donne $k$ entiers $x_1, x_2, \ldots , x_k$ strictement positifs ($k$ non nul), et on définit $n$ par :

\[n + 1  = \left(1 + x_1\right)\left(1 + x_2\right)\ldots \left(1 + x_k\right).\]

Montrer qu’il existe des entiers $d_1, d_2, \ldots, d_k$ uniques avec $1 \leqslant d_1 < d_2 < \ldots < d_k$
tels que la boîte de poids composée de $x_1$ poids de $d_1$ grammes, $x_2$ poids de $d_2$
grammes,\ldots , $x_k$ poids de $d_k$ grammes soit parfaite.
\item Déterminer toutes les boîtes de poids parfaites de masse totale \np{1993}~grammes.
\end{enumerate}

\begin{center}\textbf{{\color{blue} Exercice 2}}
\end{center}

Soit $n$ un entier strictement positif donné.

\medskip

\begin{enumerate}
\item Existe-t-il $2n+ 1$ entiers naturels consécutifs $a_0, a_1, \ldots , a_{2n}$, rangés dans l'ordre croissant, vérifiant :

\[a_0 + a_1 + \ldots + a_n = a_{n+1} + \ldots + a_{2n} ?\]

\item Existe-t-il $2n+ 1$ entiers naturels consécutifs $a_0, a_1, \ldots , a_{2n}$, rangés dans l'ordre croissant, vérifiant :

\[a_0^2 + a_1^3 + \ldots + a_n^2 = a_{n+1}^2 + \ldots + a_{2n}^2 ?\]

\item Existe-t-il $2n+ 1$ entiers naturels consécutifs $a_0, a_1, \ldots , a_{2n}$, rangés dans l'ordre croissant, vérifiant :

\[a_0^3 + a_1^3 + \ldots + a_n^3 = a_{n+1}^3 + \ldots + a_{2n}^3 ?\]

Pour cette question, on pourra étudier les variations de la fonction $f$ définie sur $\R$ par:

\[f(x) = (x - n)^3+(x- n + 1)^3+ \ldots + x^3- (x+1)^3 - \ldots  - (x + n)^3.\]

On montrera que l'équation $f(x) = 0$ admet une solution unique $x_n$ vérifiant : 

\[3n(n +1) < x_n < 3n(n +1) +1.\]

On pourra utiliser l'égalité : $1^3 + 2^3 + \ldots + n^3 = \left[\dfrac{n(n+1)}{2}\right]^2.$
\end{enumerate}

\begin{center}\textbf{{\color{blue} Exercice 3}}
\end{center}

Soit $f$ une application de l'ensemble $\Z$ des entiers relatifs dans l'ensemble $\R$ des réels.

On suppose que $f$ est minorée et vérifie : pour tout entier relatif $n$,

\[f(n) > \dfrac12[f(n + 1) + f(n - 1)].\]

Montrer que l'application $f$ est constante.

\begin{center}\textbf{{\color{blue} Exercice 4}}
\end{center}

Soit dans le plan un disque $D$ de rayon 1.

\medskip

\begin{enumerate}
\item Montrer qu'il est impossible de recouvrir le disque $D$ avec deux disques de même
rayon $r$, lorsque $r$ est strictement inférieur à 1.
\item Montrer que, pour certaines valeurs de $r,\: r < 1$, il est possible de recouvrir le disque $D$ avec trois disques de même rayon $r$. 

Quelle est la plus petite valeur de $r$ permettant un tel recouvrement ?
\end{enumerate}

\begin{center}\textbf{{\color{blue} Exercice 5}}
\end{center}

\begin{enumerate}
\item Soit A et B deux points distincts de l'espace.
	\begin{enumerate}
		\item Parmi les triangles MAB d'aire donnée, quels sont ceux de périmètre minimal ?
		\item Parmi les triangles MAB de périmètre donné, quels sont ceux d'aire maximale ?
	\end{enumerate}	
\item Soit, dans un tétraèdre de volume $V$,\: $a, b, c, d$ les longueurs de quatre arêtes, telles que trois quelconques d'entre elles ne soient pas coplanaires, et $L = a + b + c + d$.

Déterminer la valeur maximale du quotient $\dfrac{V}{L^3}$.
\end{enumerate}
\end{document}